% 2009 年广东卷（理科）第 8 题：甲、乙两车的速度—时间曲线。
% 来源 PDF：.work/ocr-local/sources/2009广东理.pdf
% SHA-256：5256b35235a042e254d36c111af3eda0f5dd8f552664021d0bace5c05582ad04
% 原页：第 1 页（A3 四栏）第 2 栏，印刷页 519。
% 绘制依据：v(t) 纵轴、t 横轴；v甲 自原点陡升后上凸（先快后慢），v乙 自原点平缓
%   后下凸（先慢后快），两线在 t0 之前相交；t0、t1 两处画竖直虚线至 v甲 曲线。
\begin{tikzpicture}[scale=1.0, semithick, every node/.style={inner sep=2pt}]
  \draw[->,>=stealth] (0,0) -- (5.6,0) node[below] {$t$};
  \draw[->,>=stealth] (0,0) -- (0,3.5) node[right] {$v(t)$};
  \node[below left] at (0,0) {$O$};

  % v甲：陡升后趋缓
  \draw[thick] (0,0) .. controls (0.55,1.35) and (1.5,2.05) .. (2.6,2.25)
    .. controls (3.4,2.4) and (4.1,2.6) .. (4.7,3.0);
  % v乙：平缓起步后急速上升
  \draw[thick] (0,0) .. controls (1.6,0.18) and (2.6,0.75) .. (3.35,1.75)
    .. controls (3.85,2.35) and (4.3,2.9) .. (4.6,3.35);

  \draw[dashed] (3.55,0) -- (3.55,2.44);
  \draw[dashed] (4.0,0) -- (4.0,2.56);
  \node[below] at (3.55,0) {$t_0$};
  \node[below] at (4.0,0) {$t_1$};

  \node at (1.95,2.55) {$v_{\text{甲}}$};
  \node at (1.75,0.8) {$v_{\text{乙}}$};
\end{tikzpicture}
